  %%%%%%%%%%%%%%%%%%%%%%%%%%%%%%%%%%%%%%%%%%%%%%%%%%%%%%%%%%%
\documentclass[12pt]{book}
\usepackage{fancyhdr}
\usepackage{makeidx} % para criar indice remissivo
\usepackage[utf8]{inputenc}
\usepackage[brazil]{babel}
\usepackage{hyperref}
\hypersetup{
colorlinks   = true, %Colours links instead of ugly boxes
  urlcolor     = cyan, %Colour for external hyperlinks
  linkcolor    = red, %Colour of internal links
  citecolor   = blue %Colour of citations
}
\usepackage[notcite,notref]{showkeys}
\usepackage{amsfonts}
\usepackage{mathtools}
\usepackage{amsthm,amssymb}
\usepackage{graphicx}
\usepackage{epsfig}
\usepackage{color}
\usepackage{enumerate}
\usepackage{wrapfig}

\usepackage{xspace,setspace}
\usepackage[rm]{titlesec}
\usepackage{fancybox}
\usepackage[hmargin=2cm,vmargin=2.5cm]{geometry}
\usepackage{tikz}
\usepackage{tcolorbox}

\usepackage[nodayofweek]{datetime}

\renewcommand{\baselinestretch}{1.2}
\allowdisplaybreaks

\usepackage[
type={CC},
modifier={by-nc},
version={4.0},
lang={brazilian},
]{doclicense}

 %%%%%%%%%%%%%%%%%%%%%%%%%%%%%%%%%%%%%%%%%%%%%%%%%%%%%%%%%%%%%%

% ----- teorema box ---------------------------
% \usepackage[framemethod=TikZ]{mdframed}
% \newcounter{theo}[section]\setcounter{theo}{0}
% \renewcommand{\thetheo}{\arabic{section}.\arabic{theo}}


% \newenvironment{proposition}[2][]{%
% 	\refstepcounter{theo}
	

% \ifstrempty{#1}%
% % if condition (without title)
% {\mdfsetup{%
% 		frametitle={%
% 			\tikz[baseline=(current bounding box.east),outer sep=0pt]
% 			\node[anchor=east,rectangle,fill=blue!20]
% 			{\strut Proposição~\thetheo};}
% 	}%
% 	% else condition (with title)
% }{\mdfsetup{%
% 		frametitle={%
% 			\tikz[baseline=(current bounding box.east),outer sep=0pt]
% 			\node[anchor=east,rectangle,fill=blue!20]
% 			{\strut Theorem~\thetheo:~#1};}%
% 	}%
% }%
% % Both conditions
% \mdfsetup{%
% 	innertopmargin=10pt,linecolor=blue!20,%
% 	linewidth=2pt,topline=true,%
% 	frametitleaboveskip=\dimexpr-\ht\strutbox\relax%
% }
	
% 	\begin{mdframed}[]\relax}{%
% \end{mdframed}}


% \newenvironment{definition}[2][]{%
% 	\refstepcounter{theo}
	
	
% 	\ifstrempty{#1}%
% 	% if condition (without title)
% 	{\mdfsetup{%
% 			frametitle={%
% 				\tikz[baseline=(current bounding box.east),outer sep=0pt]
% 				\node[anchor=east,rectangle,fill=blue!20]
% 				{\strut Definição~\thetheo};}
% 		}%
% 		% else condition (with title)
% 	}{\mdfsetup{%
% 			frametitle={%
% 				\tikz[baseline=(current bounding box.east),outer sep=0pt]
% 				\node[anchor=east,rectangle,fill=blue!20]
% 				{\strut Theorem~\thetheo:~#1};}%
% 		}%
% 	}%
% 	% Both conditions
% 	\mdfsetup{%
% 		innertopmargin=10pt,linecolor=blue!20,%
% 		linewidth=2pt,topline=true,%
% 		frametitleaboveskip=\dimexpr-\ht\strutbox\relax%
% 	}
	
% 	\begin{mdframed}[]\relax}{%
% \end{mdframed}}


% -----------------------------------------------------------




%%%%%%%%%%%%%%%%%%%%%%%%%%%%%%%%%%%%%%%%%%%%%%%%%%%%%%%%%%%%%%
\nonstopmode

\newtheorem{theorem}{Teorema}[chapter]
\newtheorem{lemma}[theorem]{Lema}
\newtheorem{corollary}[theorem]{Corolário}
\newtheorem{proposition}[theorem]{Proposição}

\theoremstyle{definition}
\newtheorem{definition}[theorem]{Definição}
\newtheorem{example}[theorem]{Exemplo}
\newtheorem{xca}[theorem]{Exercício}

\theoremstyle{remark}
\newtheorem{remark}[theorem]{Observação}

\numberwithin{section}{chapter}
\numberwithin{equation}{chapter}


%%%%%%%%%%%%%%%%%%%%%%%%%%%%%%%%%%%%%%%%%%%%%%%%%%%%%%%%%%%%%%
\newcommand{\R}{\mathbb{R}}				      % conjunto dos reais
\newcommand{\C}{\mathbb{C}}	
\renewcommand{\Re}{\operatorname{Re}}	
\newcommand{\dt}[1]{{\color{blue}\textbf{#1}}}
\newcommand{\ep}{\varepsilon}
\newcommand{\id}{\operatorname{I}}
\newcommand{\nl}[1]{\|#1\|_{\mathcal{L}(X)}}
\newcommand{\rla}{R(\lambda,A)}

%%%%%%%%%%%%%%%%%%%%%%%%%%%%%%%%%%%%%%%%%%%%%%%%%%%%%%%%%%%%%%%

%----------------------  Caixa para Instruções  --------------------------------


%---------------------------------------------------------------------------


%-----------------------------------------------------------------------
%                         ESTILO DO INÍCIO DO CAPÍTULO
%-----------------------------------------------------------------------
\font\largefont= pzcmi scaled 3500 \font\numberfont= pzcmi scaled
3000
\titleformat{\chapter}[display]
   {\normalfont\Large % roman
    }
  { \filright
   \rule[32pt]{.7\linewidth}{4pt}
   \hspace{-8pt}
      \shadowbox{
   \begin{minipage}{.15\linewidth}
     \begin{center}
          \textsl{\large \chaptertitlename}\\
       \vspace{1ex}
       {\numberfont \thechapter}\\
       \vspace{1ex}
     \end{center}
   \end{minipage}}}
  {-10pt}
  {\filcenter  \sl  \Huge }
  [\vspace{-1cm}\singlespacing\hfill\rule{.8\textwidth}{0.5pt}\\
     \vskip-2.8ex\hfill\rule{.7\textwidth}{4pt}\onehalfspacing]
\titlespacing{\chapter}{0pt}{*4}{*1}
%%%%%%%%%%%%%%%%%%%%%%%%%%%%%%%%%%%%%%%%%%%%%%%%%%%%%%%%%%%%
\makeindex

\begin{document}
\begin{titlepage}
\begin{center}
\includegraphics[scale=0.15]{brasao-cor.png}~\\
UNIVERSIDADE FEDERAL FLUMINENSE \\
PÓS-GRADUAÇÃO DE MATEMÁTICA DA UFF\\[4cm]
{\Huge \textbf{Semigrupos de Operadores Lineares}}\\
{\Large
\begin{equation*}
\begin{cases}
x'(t)={\color{red}A}x+{\color{black}f},& t\in [0,+\infty)\\
x(0)=x_0
\end{cases}
\end{equation*}
\begin{equation*}
x(t)=e^{t\color{red}A}x_0+\int_0^t e^{(t-s){\color{red}A}}f(s)ds
\end{equation*}}\\[1cm]
\textbf{Luiz Viana}\footnote{Departamento de Análise, Universidade Federal Fluminense, Niterói, RJ.}\\
%{\footnotesize Departamento de Análise, Universidade Federal Fluminense, Niterói, RJ.}\\[0.5cm]
\textbf{Reginaldo Demarque}\footnote{Departamento de Ciências da Natureza, Universidade Federal Fluminense, Rio das Ostras, RJ.}\\
% {\footnotesize Departamento de Ciências da Natureza, Universidade Federal Fluminense, Rio das Ostras, RJ.}
\vfill
% Bottom of the page

{\large Niterói, 2024}

% Un-adjusted
[Compilado \today, \currenttime] 
\par
\end{center}
\end{titlepage}


\null\vfill
\noindent
\begin{minipage}{0.8\textwidth}
Semigrupos de Operadores Lineares \copyright 2025 by Reginaldo Demarque and Luiz Viana tem a licença \ \ \doclicenseNameRef \ \ \doclicenseIcon. Para ver uma cópia da licença, visite \url{\doclicenseURL} 
\end{minipage}
\doclicenseImage[imagewidth=7em] 


\newpage






%%%%%%%%%%%%%%%%%%%%%%%%%%%%%%%%%%%%%%%%%%%%%%%%%%%%%%%%%%%%
%%%%%
%%%%%       PÁGINAS PRÉ-TEXTUAIS
%%%%%
%%%%%%%%%%%%%%%%%%%%%%%%%%%%%%%%%%%%%%%%%%%%%%%%%%%%%%%%%%%%
\frontmatter
\tableofcontents
%----------------------------------------------------------------------
%                         ESTILO DA PÁGINA
%----------------------------------------------------------------------

\pagestyle{fancy}
\fancyhead[LO]{\slshape \rightmark}
\fancyhead[RE]{\slshape \leftmark}
\fancyhead[LE,RO]{\thepage}
\fancyfoot{}
\renewcommand{\headrulewidth}{0.4pt}
\renewcommand{\chaptermark}[1]{%
 \markboth{\chaptername
\ \thechapter.\ #1}{}}
\renewcommand{\sectionmark}[1]{\markright{\thesection.\ #1}}
%%%%%%%%%%%%%%%%%%%%%%%%%%%%%%%%%%%%%%%%%%%%%%%%%%%%%%%%%%%%%%%%%%%%%%%%%%%%%%%%%%

%%%%%%%%%%%%%%%%%%%%%%%%%%%%%%%%%%%%%%%%%%%%%%%%%
%%%%
%%%% Capitulos da Monografia
%%%%
%%%%%%%%%%%%%%%%%%%%%%%%%%%%%%%%%%%%%%%%%%%%%%%%


\mainmatter
%\chapter{Preliminares}

\chapter{Noções Básicas sobre Espaços de Banach}

\section{Espaços Normados}

No decorrer do presente capítulo, denotaremos por $\mathbb K$ o corpo dos números reais ou o corpo dos números complexos. Além disso, quando $(M,d)$ for um espaço métrico, cada bola aberta, cada bola fechada e cada esfera em $M$ serão denotadas por
\[
\displaystyle B(a;r)=\{y\in M; d(y,a)<r\},
\]

\[
\displaystyle B(a;r)=\{y\in M; d(y,a)\leq r\}
\]
e
\[
\displaystyle B(a;r)=\{y\in M; d(y,a)=r\},
\]
respectivamente, onde $a\in M$ e $r>0$.

\begin{definition}\label{norma}
    Seja $X$ um $\mathbb K$-espaço vetorial. Uma aplicação
    \[
    \displaystyle \| \cdot \|: X\longrightarrow \R
    \]
    é dita uma \dt{norma} \index{norma} se, para quaisquer $x,y\in X$ e $\lambda \in \mathbb K$, as seguintes condições se verificarem:
    \begin{itemize}
    \item[(a)] $\|x\|\geq 0$;
    \item[(b)] Se $\|x\|=0$, então $x=0$;
    \item[(c)] $\|\lambda x\|=|\lambda| \|x\|$;
    \item[(d)] $\|x+y\|\leq \|x\|+\|y\|$.
    \end{itemize}
    Nesse caso, o par $(X,\|\cdot\|)$ é dito um \dt{espaço normado} \index{Espaço!normado}.
\end{definition}

\begin{remark}
Em um espaço normado $(X,\|\cdot\|)$, valem:
    \begin{itemize}
    \item[(a)] $\|0\|=0$;
    \item[(b)] $|\|x\|-\|y\|| \leq \|x-y\|$ para quaisquer $x,y \in X$.
    \end{itemize}
\end{remark}

\begin{example}\label{ex21}
    Dado um número inteiro positivo $n$, não é difícil verificar que
    \[
    \displaystyle \|\cdot \|_0 : (x_1,\cdots , x_n) \in \mathbb K^n \longmapsto \bigg( \sum_{j=1}^{n} |x_j|^2\bigg)^{1/2} \in \R ,
    \]
    \[
    \displaystyle \|\cdot \|_1 : (x_1,\cdots , x_n) \in \mathbb K^n \longmapsto \max_{1\leq j \leq n} |x_j| \in \R 
    \]
    e
    \[
    \displaystyle \|\cdot \|_2 : (x_1,\cdots , x_n) \in \mathbb K^n \longmapsto \sum_{j=1}^{n} |x_j| \in \R 
    \]
    são normas e $\mathbb K ^n$.
\end{example}

\begin{example}\label{limitadas}
Seja $A$ um conjunto não vazio. Denotemos por $\mathcal B (A)$ o conjunto de todas as funções limitadas $f: A\longrightarrow \mathbb K$. Dados $f,g \in \mathcal B (A)$ e $\lambda \in \mathbb K$, definamos
\begin{itemize}
\item $(f+g)(x):=f(x)+g(x)$; 
\item $(\lambda f)(x) = \lambda f(x)$,
para cada $x\in A$.
\end{itemize}
Com as operações operações de adição e multiplicação por escalar, pontualmente dadas acima, $\mathcal B (A)$ é um $\mathbb K$-espaço vetorial. Além disso, não é difícil constatar que
\[
\displaystyle f\in \mathcal B (A) \longmapsto \| f\| = \sup_{x\in A} |f(x)| \in \R
\]
é uma norma em $\mathcal B (A)$. Observemos ainda que a convergência de sequências em $\mathcal B (A)$ é exatamente a noção de convergência uniforme. 
\end{example}

\begin{definition}
    Com as notações do Exemplo \ref{limitadas},
    \[
    \ell_{\infty} := \mathcal B (\mathbb N),
    \]
    que é o espaço normado de todas as sequências limitadas cujos termos pertencem a $\mathbb K$.
\end{definition}

\begin{example}\label{lp}
    Seja $p\in [1,+\infty )$ e definamos 
    \[
    \displaystyle \ell_p =\left\{ (x_n )_{n=1}^{\infty}; \sum_{n=1}^{\infty} |x_n|^p <+\infty \right\}
    \]
    
    Mostraremos que as operações 
    \[
    \bigg( (x_n )_{n=1}^{\infty},(y_n )_{n=1}^{\infty} \bigg) \in \ell_p \times \ell_p \longmapsto (x_n +y_n )_{n=1}^{\infty} \in \ell_p
    \]
    e
    \[
    \bigg( \lambda ,(x_n )_{n=1}^{\infty} \bigg) \in \mathbb K \times \ell_p \longmapsto (\lambda x_n )_{n=1}^{\infty} \in \ell_p
    \]
    estão bem definidas e que a aplicação 
    \[
    \displaystyle (x_n )_{n=1}^{\infty} \in \ell_p \longmapsto \bigg( \sum_{n=1}^{\infty} |x_n|^p \bigg)^{1/p} \in \R .
    \]
    Tais objetivos serão alcançados no Corolário \ref{dms}.
\end{example}

\begin{proposition}[Desigualdade de Hölder para somas] \label{dhsomas}
    Sejam $p,q \in (1,+\infty)$ tais que 
    \[
    \displaystyle \frac{1}{p} + \frac{1}{q} =1.
    \]
    Se $(x_1, \cdots , x_n) \in \mathbb K^n$ e $(y_1, \cdots , y_n) \in \mathbb K^n$, então 
    \[
    \displaystyle \sum_{j=1}^{n} |x_j y_j|\leq \bigg( \sum_{j=1}^{n} |x_j|^p \bigg)^{1/p} \bigg( \sum_{j=1}^{n} |y_j|^q \bigg)^{1/q} .
    \]
\end{proposition}
\begin{proof}
    Em virtude do Teorema do Valor Médio, não é difícil constatar que, para quaisquer $a,b\in [0,+\infty )$ e $\alpha \in [0,1]$, tem-se
    \begin{equation}\label{TVM}
    \displaystyle a^{\alpha} b^{1-\alpha} \leq \alpha a + (1-\alpha )b.
    \end{equation}
    Como claramente podemos supor que
    \[
    \sum_{j=1}^{n} |x_j|^p >0 \text{  e  } \sum_{j=1}^{n} |y_j|^p >0,
    \]
    consideremos 
    \[
    \displaystyle a_{m} = \frac{|x_{m}|^p}{\sum_{j=1}^{n} |x_j|^p} \text{  e  } b_{m} = \frac{|y_{m}|^q}{\sum_{j=1}^{n} |y_j|^p},
    \]
    para cada $m\in \{1,\cdots , n\}$. Logo, tomando $\alpha = \frac{1}{p}$ e aplicando \eqref{TVM}, temos
    \[
    \displaystyle a_{m}^{\alpha} b_{m}^{1-\alpha} \leq \alpha a_m + (1-\alpha )b_m ,
    \]
    isto é,
    \begin{equation}\label{ineq}
    \displaystyle \frac{|x_{m}|}{\bigg(\sum_{j=1}^{n} |x_j|^p \bigg)^{1/p}} \cdot \frac{|y_{m}|}{\bigg(\sum_{j=1}^{n} |y_j|^p \bigg)^{1/q}}
    \leq \frac{1}{p} \cdot \frac{|x_{m}|^p}{\sum_{j=1}^{n} |x_j|^p}+ \frac{1}{q} \cdot \frac{|y_{m}|^q}{\sum_{j=1}^{n} |y_j|^p},
    \end{equation}
    para cada $m\in \{ 1,\cdots , n\}$. Somando membro a membro as $n$ relações descritas em \eqref{ineq}, fica demonstrada a desigualdade de Hölder do enunciado.
\end{proof}

\begin{corollary}[Desigualdade de Hölder para séries]
    Sejam $p,q \in (1,+\infty)$ satisfazendo $\frac{1}{p} + \frac{1}{q} =1$. Se $(x_n)_{n=1}^{\infty} \in \ell _p$ e $(y_n)_{n=1}^{\infty} \in \ell _q$, então $(x_n y_n)_{n=1}^{\infty} \in \ell _1$ e vale
    \[
    \displaystyle \sum_{n=1}^{\infty} |x_n y_n|\leq \bigg( \sum_{n=1}^{\infty} |x_n|^p \bigg)^{1/p} \bigg( \sum_{n=1}^{\infty} |y_n|^q \bigg)^{1/q} .
    \]
\end{corollary}
\begin{proof}
Segue imediatamente da Proposição \ref{dhsomas}.
\end{proof}


\begin{proposition}[Desigualdade de Minkowski para somas]
Seja $p\in (1,+\infty)$. Se $(x_1, \cdots , x_n) \in \mathbb K^n$ e $(y_1, \cdots , y_n) \in \mathbb K^n$, então 
    \[
    \displaystyle \bigg( \sum_{j=1}^{n} |x_j + y_j|^p \bigg)^{1/p} \leq \bigg( \sum_{j=1}^{n} |x_j|^p \bigg)^{1/p} +\bigg( \sum_{j=1}^{n} |y_j|^p \bigg)^{1/p} .
    \]
\end{proposition}

\begin{proof}
Como claramente podemos supor que
\[
\displaystyle \sum_{j=1}^{n} |x_j + y_j|^p >0,
\]
a Proposição \ref{dhsomas} nos dá

\begin{align}
\displaystyle \sum_{j=1}^{n} |x_j + y_j|^p
&=\sum_{j=1}^{n} |x_j + y_j|^{p-1} |x_j + y_j| \nonumber \\
&\leq \sum_{j=1}^{n} |x_j + y_j|^{p-1} |x_j|
+\sum_{j=1}^{n} |x_j + y_j|^{p-1} |y_j| \nonumber \\
&\leq \bigg( \sum_{j=1}^{n} |x_j + y_j|^{(p-1)q} \bigg)^{1/q} \bigg( \sum_{j=1}^{n} |x_j|^p \bigg)^{1/p} \nonumber\\
&\hspace{1cm}+\bigg( \sum_{j=1}^{n} |x_j + y_j|^{(p-1)q} \bigg)^{1/q} \bigg( \sum_{j=1}^{n} |y_j|^p \bigg)^{1/p} \nonumber \\
&=\bigg( \sum_{j=1}^{n} |x_j + y_j|^{p} \bigg)^{1/q} \bigg[ \bigg( \sum_{j=1}^{n} |x_j|^p \bigg)^{1/p} + \bigg( \sum_{j=1}^{n} |y_j|^p \bigg)^{1/p}\bigg] \label{ineqm},
\end{align}
onde $q=\frac{p}{p-1}$. Multiplicando os membros de \eqref{ineqm} por $\bigg( \sum_{j=1}^{n} |x_j + y_j|^p \bigg)^{-1/q}$, fica demonstrada a desigualdade de Minkowski do enunciado.
\end{proof}

\begin{corollary}[Desigualdade de Minkowski para séries]\label{dms}
    Seja $p\in (1,+\infty)$. Se $x=(x_n)_{n=1}^{\infty} \in \ell_p $ e $y=(y_n)_{n=1}^{\infty} \in \ell_q$, então $(x_n + y_n)_{n=1}^{\infty} \in \ell_p$ e vale
     \[
    \displaystyle \bigg( \sum_{n=1}^{\infty} |x_n + y_n|^p \bigg)^{1/p} \leq \bigg( \sum_{n=1}^{\infty} |x_n|^p \bigg)^{1/p} +\bigg( \sum_{n=1}^{\infty} |y_n|^p \bigg)^{1/p} .
    \]
    Em outras palavras, $x+y \in \ell_p$, ou seja, a adição em $\ell_p$ apresentada no Exemplo \ref{lp} encontra-se bem definida e, além disso, vale a desigualdade triangular
    \[
    \displaystyle \|x+y\|_p \leq \|x\|_p +\|y\|_p.
    \]
    Assim, sendo imediata a boa definição da multiplicação por escalar em $\ell_p$, também declarada no Exemplo \ref{lp}, resulta que $(\ell_p , \|\cdot\|_p)$ é um espaço vetorial normado.
\end{corollary}    

\begin{example}\label{ex23}
    Sendo $M$ um espaço métrico, denotemos por $\mathcal C_b (M)$ o subconjunto de $\mathcal B (M)$ formado por todas as funções contínuas e limitadas de $M$ em $\mathbb K$. Não é difícil constatar que $\mathcal C_b (M)$ é um subespaço vetorial fechado de $\mathcal B (M)$.
\end{example}

\begin{example}
Denotemos por
\begin{itemize}
\item $\textbf{c}$ o conjunto de todas as sequências convergentes em $\mathbb K$;
\item $\textbf{c}_0$ o conjunto de todas as sequências convergentes em $\mathbb K$ convergindo para zero;
\item $\textbf{c}_{00}$ o conjunto de todas as sequências $(x_n)_{n=1}^{\infty}$ em $\mathbb K$ com a seguinte propriedade: existe um inteiro positivo $n_0$ tal que $x_n=0$ para todo $n\geq n_0$.
\end{itemize}
Claramente, 
\[
\textbf{c}_{00} \subset \textbf{c}_{0} \subset \textbf{c} \subset \ell_{\infty}.
\]
Além disso, não é difícil constatar que $\textbf{c}$ é um subespaço vetorial fechado de $\ell_{\infty}$.
\end{example}

O próximo resultado traz importantes caracterizações das aplicações lineares contínuas.

\begin{proposition}\label{continuous}
Sejam $X$ e $Y$ dois espaços normados, e consideremos uma aplicação linear $T:X\longrightarrow Y$. As seguintes condições são equivalentes:

\begin{itemize}
\item[(a)] $T$ é uniformemente contínua;
\item[(b)] $T$ é contínua;
\item[(c)] $T$ é contínua em $0\in X$;
\item[(d)] Existe $C>0$ tal que $\|T(x)\|_Y \leq C \|x\|_X$ para todo $x\in X$.
\end{itemize}
\end{proposition}

\begin{proof}
É claro que $(a)\Longrightarrow (b) \Longrightarrow (c)$. Para vermos que $(c)\Longrightarrow (d)$, , tomemos $\varepsilon >0$. Como $T$ é contínua em $0\in X$, existe $\delta >0$ tal que
\[
\displaystyle \|T(x)\|_Y =\|T(x) - T(0)\|_Y <\varepsilon =1,
\]
sempre que $x\in X$ e $\|x\|_X = \|x-0\|_X <\delta$. Assim, para todo $w\in X\setminus \{0\}$, temos
\[
\displaystyle \bigg\|T\bigg( \frac{\delta}{2} \bigg( \frac{w}{\|w\|_X} \bigg)\bigg) \bigg\|_Y <1.
\]
Logo, 
\[
\|T(w)\|_Y \leq \frac{2}{\delta}  \|w\|_X \text{  para todo  } w\in X,
\]
inclusive se $w=0$.

Agora, para vermos que $(d)\Longrightarrow (a)$, basta observarmos que
\[
\|T(x) - T(y)\|_Y =\|T(x-y)\|_Y \leq C \|x-y\|_X
\]
para quaisquer $x,y\in X$. Ou seja: $T$ é uma aplicação lipschitziana e, portanto, uniformentente contínua.
\end{proof}

\begin{definition}
    Sejam $X$ e $Y$ dois espaços normados. Denotaremos por $\mathcal L (X,Y)$ o espaço vetorial de todas as transformações lineares de $X$ em $Y$, com as operações de adição e multiplicação por escalar definidas pontualmente. Cabe observar que:
    \begin{itemize}
    \item[(a)] Quando $X=Y$, escreveremos $\mathcal L (X)$ em vez de $\mathcal L (X,X)$;
    \item[(b)] Quando $Y=\mathbb K$, denotaremos $\mathcal L (X,\mathbb K)$ por $X'$, que é conhecido como \textbf{o dual topológico} de X;
    \item[(c)] O conjunto de todos os funcionais lineares $\varphi : X\longrightarrow \mathbb K$, contínuos ou não, será denotador por $X^{\ast}$, que é conhecido como \textbf{o dual algébrico} de X.
    \end{itemize}
\end{definition}

\begin{remark}
Seja $X$ um espaço normado. É sempre verdade que que $X' \subset X^{\ast}$. Além disso, quando $X$ tem dimensão infinita, sempre temos $X'\neq X^{\ast}$. De fato, em virtude do Lema de Zorn, sabemos que $X$ possui uma base (de Hamel)
\[
\displaystyle \mathcal B = \{ x_i; i\in I\},
\]
onde $I$ é um conjunto infinito. Podemos supor que $\| x_i \|_X =1$ para todo $i\in I$. Agora, seja 
\[
\displaystyle J = \{ i_1, i_2, \cdots , i_k, \cdots \}
\]
um subconjunto infinito e enumerável de $I$. Logo, existe um único funcional linear $\varphi : X \longrightarrow \mathbb K$ tal que 
\begin{itemize}
\item $\varphi (x_{i_k})=k$ para cada $k\in \mathbb N$;
\item $\varphi (x_i)=0$ se $i\in I\setminus J$.
\end{itemize}
Um vez que 
\[
\displaystyle \mathbb N \subset \left\{ \frac{|\varphi (x)|}{\|x\|_{X}}; x\in X\setminus \{0\}\right\},
\]
não existe $C>0$ de modo que valha $|\varphi (x)|\leq \|x\|_X$ para todo $x\in X$. Pela Proposição \ref{continuous}, $\varphi$ não é contínua, isto é, $\varphi \in X^{\ast} \setminus X'$.
\end{remark}

\begin{remark}
Sejam $X$ e $Y$ dois espaços normados. Não é difícil constatar que:
\begin{itemize}
\item[(a)] Se $X$ tem dimensão finita, então toda aplicação linear $T:X\longrightarrow Y$ é contínua;
\item[(b)] Se X tem dimensão infinita, sempre existe uma aplicação linear $T:X\longrightarrow Y$ que não é contínua.
\end{itemize}
\end{remark}

\begin{definition}
Sejam $X$ e $Y$ dois espaços normados. Uma aplicação linear $T:X \longrightarrow Y$ é dita \textbf{limitada} se 
\[
\displaystyle \sup_{x\in X\setminus \{0\}} \frac{\|T(x)\|_Y}{\|x\|_X} <\infty.
\]
\end{definition}

\begin{remark}
Sejam $X$ e $Y$ dois espaços normados. Em virtude da Proposição \ref{continuous}, uma aplicação linear $T:X\longrightarrow Y$ é contínua se, e somente se, é limitada. 
\end{remark}

\begin{example}
Sejam $X$ e $Y$ dois espaços normados. Vejamos que 
\[
T\in \mathcal L (X,Y) \longmapsto \|T\|_{\mathcal L (X,Y)} = \sup_{x\in X\setminus \{0\}} \frac{\|T(x)\|_Y}{\|x\|_X} \in \R
\]
é uma norma em $\mathcal L (X,Y)$. De fato, as condições (a), (b) e (c) da Definição \ref{norma} são claras. Além disso, se $T,S\in \mathcal L (X,Y)$, temos

\begin{equation*}
\displaystyle \|(T+S)(x)\|_Y=\|T(x)+S(x)\|_Y \leq \|T(x)\|_Y +\|S(x)\|_Y \leq (\|T\|_{\mathcal L (X,Y)} + \|S\|_{\mathcal L (X,Y)})\|x\|_X
\end{equation*}
para todo $x\in E$, ou seja,
\begin{equation*}
    \displaystyle \|T+S\|_{\mathcal L (X,Y)}=\sup_{x\in X\setminus \{0\}} \frac{\|(T+S)(x)\|_Y}{\|x\|_X} \leq \|T\|_{\mathcal L (X,Y)}+\|S\|_{\mathcal L (X,Y)}.
\end{equation*}
Assim, temos realmente uma norma em $\mathcal L (X,Y)$.
\end{example}

\begin{definition}
Sejam $X$ e $Y$ dois espaços normados. Uma aplicação linear $T:X\longrightarrow Y$ é dita um \textbf{isomorfismo topológico} se for um homeomorfismo. 
\end{definition}

\begin{definition}
Sejam $\|\cdot\|_1$ e $\|\cdot\|_2$ duas normas em um mesmo espaço vetorial $X$. Dizemos que tais normas são textbf{equivalentes} se a aplicação identidade
\[
\displaystyle I_X : (X,\|\cdot\|_1) \longrightarrow (X,\|\cdot\|_2)
\]
for um isomorfismo topológico.
\end{definition}

\begin{corollary}
Sejam $\|\cdot\|_1$ e $\|\cdot\|_2$ duas normas em um mesmo espaço vetorial $X$. As seguintes condições são equivalentes:
\begin{itemize}
\item[(a)] $\|\cdot\|_1$ e $\|\cdot\|_2$ são equivalentes;
\item[(b)] Existem constantes $A>0$ e $B>0$ tais que
\[
A\|x\|_2 \leq \|x\|_1 \leq \|x\|_2
\]
para todo $x\in X$.
\end{itemize}
\end{corollary}

\begin{proof}
Basta aplicar a Proposição \ref{continuous}.
\end{proof}

\section{Espaços de Banach}

\begin{definition}
    Seja $X$ um espaço normado. Dizemos que $X$ é um \textbf{espaço de Banach} se ele for um espaço métrico completo com a métrica proveniente da sua norma.
\end{definition}

\begin{remark}\label{subespaço}
Se $X$ é um espaço métrico, então cada subespaço fechado $M$ de $X$ é também um espaço de Banach, com a norma induzida pela norma de $X$.
\end{remark}

\begin{example}
Dado um inteiro positivo $n$, $\mathbb K ^n$ é um espaço de Banach com a norma $\|\cdot\|_0$ do Exemplo \ref{ex21}. Na verdade, como as normas $\| \cdot\|_0$, $\| \cdot\|_1$ e $\| \cdot\|_2$ são duas a duas equivalentes, $(\mathbb K ^n ,$\| \cdot\|_j )$ é um espaço de Banach para todo $j\in \{1,2,3\}$. 
\end{example}

\begin{example}
Dado um conjunto $A\neq \emptyset$, não  é difícil constatar que $\mathcal B (A)$, introduzido no Exemplo \ref{limitadas}, é um espaço de Banach. Em particular, $\ell _{\infty}$ é um espaço de Banach.
\end{example}

\begin{example}
    Para cada espaço métrico $M$, $\mathcal C _b (M)$, apresentado no \ref{ex23} é um espaço de Banach, uma vez que é um subespaço fechado de $\mathcal B (M)$ (veja a Observação \ref{subespaço}).
\end{example}

\section{Operadores Lineares Ilimitados}

Sejam $X$ e $Y$ dois espaços de Banach. Seja $A:D(A)\subset X\longrightarrow Y$ um operador linear,  onde $D(A)$ é um subespaço de $X$, chamado de \dt{domínio}\index{Operador Linear!domínio} de $A$.
\medskip
Dizemos que $A$ é \dt{limitado (ou contínuo)} \index{Operador Linear!limtado} se $D(A)=X$ e se existe  $C>0$ tal que 
\begin{equation}\label{ltdo}
\|Ax\|_{Y}\leq C\|x\|_X, \ \forall x\in X.
\end{equation}
Denotaremos por $\mathcal{L}(X,Y)$ o espaço de Banach dos \dt{operadores lineares limitados} com norma dada por
\[\|A\|_{_{\mathcal{L}(X,Y)}}=\sup\limits_{x\in X, x\neq 0}\frac{\|Ax\|_{_Y}}{\|x\|_{_X}},\]
por $\mathcal{L}(X)=\mathcal{L}(X,X)$ e $X'=\mathcal{L}(X,\R)$.
$A$ é dito ser \dt{ilimitado} \index{Operador Linear!ilimitado} quando não satisfazer \eqref{ltdo}. Dizemos que $A$ é \dt{densamente definito} \index{Operador Linear!desamente definido} se $\overline{D(A)}=X$.

\begin{theorem}[Banach-Steinhaus]\label{th-BS}
	Sejam $X$ um espaço de Banach e $Y$ um espaço vetorial normado. Se $\{T_i\}_{i\in I}$  é uma família (não necessariamente enumerável) em $\mathcal{L}(X,Y)$ pontualmente limitada, então $\{T_i\}_{i\in I}$ é uniformemente limitada, isto é, se 
	 \[\sup_{i\in I}\|T_ix\|< \infty,\ \forall x\in X,\] 
	 então
	 \[\sup_{i\in I}\|T_i\|_{\mathcal{L}(X,Y)}<\infty.\]
\end{theorem}

\section{Integrais Vetoriais}

\begin{definition}
Sejam $X$ um espaço de Banach e $u:[a,b]\longrightarrow X$ uma aplicação tal que, para cada $\varphi\in X'$,  a função real
\[t\in [a,b] \longmapsto \langle\varphi, u(t) \rangle_{X'X}\in \R,\]
seja integrável. Dizemos que \dt{$u$ é integrável} se existe um vetor $v\in X$ que satisfaz:
\[\langle \varphi, v\rangle_{X',X}=\int_a^b \langle\varphi, u(t) \rangle_{X'X}\,dt,\ \forall \varphi\in X'.\]
Em caso afirmativo, $v$ é único e escrevemos
\[v=\int_a^b u(t)\,dt.\]
\end{definition}

\begin{proposition}\label{KthA3.2}
Se $u:[a,b]\longrightarrow X$ é {contínua}, então $u$ é integrável. Além disso, 
\begin{enumerate}
    \item $\displaystyle\left\|\int_a^b u(t)\,dt\right\|\leq \int_a^b \|u(t)\|\,dt$
    \item Se $A\in \mathcal{L}(X,Y)$, então
    \[ A\left(\int_a^b u(t)\,dt\right)=\int_a^b A(u(t))\,dt. \]
\end{enumerate}
\end{proposition}
\begin{proof}
Veja em \cite[Theorem A3.2]{kesavan2015topics}
\end{proof}

\begin{proposition}
Se $u:[a,a+h]\longrightarrow X$ é {contínua}, então
\begin{equation}\label{ineq.VM}
\lim\limits_{h\to 0^+} \frac{1}{h}\int_a^{a+h} u(t)\,dt = u(a).
\end{equation}
\end{proposition}
\begin{proof}
A função  $f:t\in [a,a+h]\longmapsto \|u(t)-u(a)\|\in \R$ é contínua. Pelo Teorema do Valor Médio para integrais, existe $\xi_h\in [a,a+h]$ tal que 
\[
\frac{1}{h}\int_a^{a+h} f(t)\,dt=f(\xi_h).
\]
Como $a\leq \xi_h\leq a+h$, se $h\to 0^+$, então $\xi_h\to a^+$. Neste caso, 
\[
\lim_{h\to 0^+}\frac{1}{h}\int_a^{a+h} f(t)\,dt=\lim_{h\to 0^+}f(\xi_h)=f(a).
\]
Isto é, 
\begin{equation*}
\lim_{h\to 0^+}\frac{1}{h}\int_a^{a+h} \|u(t)-u(a)\|\,dt=0.
\end{equation*}
Assim, 
\begin{align*}
\lim_{h\to 0^+}\left\|\frac{1}{h}\int_a^{a+h} u(t)\,dt -u(a)\right\|& =
\lim_{h\to 0^+}\left\|\frac{1}{h}\int_a^{a+h} u(t)-u(a)\,dt\right\|\\
&\leq \lim_{h\to 0^+}\frac{1}{h}\int_a^{a+h}\left\| u(t)-u(a)\right\|\,dt
=0.
\end{align*}
O que é equivalente à identidade \eqref{ineq.VM}.
\end{proof}
\chapter{Semigrupos de Classe $C^0$}

\section{Semigrupos}



\begin{definition}
     Seja $X$ um espaço de Banach. Dizemos que a aplicação $S:[0,+\infty)\longrightarrow \mathcal{L}(X)$ é um  \dt{semigrupo de operadores limitados em X} quando:
\begin{enumerate}
    \item $S(0)=\operatorname{Id}$;
    \item $S(t+s)=S(t)S(s)$, $\forall t,s\in [0,+\infty)$;

Dizemos que $S$ é \dt{de classe $C^0$ ou fortemente contínuo} se

    \item $\lim\limits_{t\to0^+} \|(S(t)-\operatorname{Id})x\|=0$, $\forall x\in X$.

Dizemos que $S$ é \dt{uniformemente contínuo} se 

\item $\lim\limits_{t\to 0^+}\|S(t)-\operatorname{Id}\|_{\mathcal{L}(X)}=0$.
\end{enumerate}
\end{definition}

   \begin{example} São exemplos de semigrupos:
        \begin{enumerate}
            \item Sejam $X$ um espaço de Banach e  $A\in \mathcal{L}(X)$. Define-se a aplicação exponencial por
            \[e^{A}=\operatorname{Id}+\sum_{n=1}^\infty \frac{A^n}{n!}.\]
            Pode-se mostrar que esta série é absolutamente convergente e define $e^{A}\in \mathcal{L}(X)$, veja \cite[Apêndice 2]{gomes1985semigrupos}. Neste caso,  $e^{tA}:[0,+\infty)\to \mathcal{L}(X)$, quando $A\in \mathcal{L}(X)$, é um semigrupo.
            \item Seja $X=C_b(\R)$  o espaço das funções $f:\R\longrightarrow \R$ uniformemente contínuas e limitadas, com a norma do $\sup$. Então $S(t)f(s)=f(t+s)$ definie um semigrupo de classe $C^0$.
        \end{enumerate}
    \end{example}


\begin{proposition}
    Se $S$ é um semigrupo de classe $C^0$ em X, então existem $\mu\geq 0$ e $M\geq 1$ tais que 
    \begin{equation}\label{des.wT}
    \|S(t)\|_{\mathcal{L}(X)}\leq Me^{\mu t},\ \forall t\geq 0.  
    \end{equation}    
    Em particular, $\|S(t)\|_{\mathcal{L}(X)}$ é uma função limitada em todo intervalo $[0,T]$. 
\end{proposition}

\begin{proof}
Vamos aplicar o Teorema de Banach-Steinhaus (Teorema \ref{th-BS} )  à família  $\{S(t)\}_{t\geq 0}$ em $\mathcal{L}(X)$, onde $X$ é Banach. Para isso, basta mostrar que esta família é pontualmente limitada. 

De fato,  $S$ é um semigrupo de classe $C^0$, então
\[\lim\limits_{t\to 0^+}\|S(t)x-x\|=0,\ \forall x\in X.\]
Com isso, dado $x\in X$, para $\ep=1$, existe $\delta>0$ tal que, se $0\leq t\leq \delta$, então
\[\|S(t)x\|\leq \|S(t)x-x\|+\|x\|\leq 1+\|x\|=C_x.\]

Portanto, pelo Teorema de Banach-Steinhaus, $\{S(t)\}_{0\leq t\leq\delta}$ é uniformemente limitada, isto é, $\exists\, M>0$ tal que
\begin{equation}\label{eq1}
\|S(t)\|_{\mathcal{L}(X)}\leq M,\, \text{ para todo  } t\in [0,\delta]. 
\end{equation}
Além disso, $M\geq \|S(0)\|_{\mathcal{L}(X)}=\|\operatorname{Id}\|_{\mathcal{L}(X)}=1$.

Por outro lado, dado $t>\delta$, pelo algorítimo da divisão, existem $n\in\mathbb{N}$ e $r\in [0,\delta)$ tais que $t=n\delta+r$. Com isso, do item 2 da definição de Semigrupo e da desigualdade \eqref{eq1}, temos que
\begin{align*}
 \|S(t)\|_{\mathcal{L}(X)} &
 \leq \|S(n\delta +r)\|_{\mathcal{L}(X)}= \|S(n\delta)S(r)\|_{\mathcal{L}(X)}=
 \|\overbrace{S(\delta)\cdots S(\delta)}^{n \text{ vezes}}S(r)\|_{\mathcal{L}(X)}\\
& \leq 
\|S(\delta)\|^n_{\mathcal{L}(X)}\|S(r)\|_{\mathcal{L}(X)}\leq M^{n+1}. 
\end{align*}
 Note que, como $n\leq t/\delta$ e $M\geq 1$, temos que $M^n\leq M^{t/\delta}$. Da desigualdade anterior,
\[\|S(t)\|_{\mathcal{L}(X)}\leq M M^{t/\delta}=Me^{\frac{t}{\delta}\log(M)}
=Me^{\mu t},\]
onde $\mu=\frac{1}{\delta}\log(M)$.

Em particular, 
\[\|S(t)\|_{\mathcal{L}(X)}\leq Me^{\mu t}\leq Me^{\mu T},\, \forall t\in [0,T].\]
Em outras palavras, $\nl{S(t)}$ é limitada em $[0,T]$.
\end{proof}


\begin{corollary}\label{Scontinua}
Todo semigrupo de classe $C^0$ é fortemente contínuo em $[0,+\infty)$, i.e., 
para todo $x\in X$, $S(\cdot)x\in C^0([0,+\infty);X)$. Em outras palavras,
\begin{center}
$t\in [0,+\infty)\longmapsto S(\cdot)x\in  X$ é contínua.
\end{center}
\end{corollary}

\begin{proof}
    Dado $x\in X$, devemos mostrar que, para todo $t\in [0,+\infty)$,
    \[\|S(t+h)x-S(t)x\|\to 0,\, \text{ quando } h\to 0.\]
Do item (3) da definição de semigrupos, já temos que
  \[\|S(h)x-S(0)x\|=\|(S(h)-\operatorname{Id})x\|\to 0,\, \text{ quando } h\to 0.\]
Isto é, $S(\cdot)x$ é contínua em $t=0$. Dado $t\in (0,+\infty)$, se $h>0$, então
\begin{align*}
    \|S(t+h)x-S(t)x\|& =\|S(t)S(h)x-S(t)x\|=\|S(t)(S(h)-\operatorname{Id})x\|\\
    & \leq \|S(t)\|_{\mathcal{L}(X)}\|(S(h)-\operatorname{Id})x\|\\
    & \leq Me^{\mu t}\|(S(h)-\operatorname{Id})x\|\to 0,\, \text{ quando } h\to 0^+.
\end{align*}

Se $h<0$, então seja $k=-h>0$. Daí, se $h\to 0^+$, então $k\to0^+$. Com isso,
\begin{align*}
    \|S(t+h)x-S(t)x\|& = \|S(t-k)x-S(t)x\|=\|S(t-k)x-S(t-k+k)x\|\\
    & =\|S(t-k)x-S(t-k)S(k)x\|=\|S(t-k)(\id-S(k)x\|\\
    & \leq \|S(t-k)\|_{\mathcal{L}(X)}\|(S(k)-\operatorname{Id})x\|\\
    & \leq Me^{\mu t}\|(S(k)-\operatorname{Id})x\|\to 0,\, \text{ quando } k\to 0^+.
\end{align*}
\end{proof}

Vamos melhorar a estimativa \eqref{des.wT} através do seguinte teorema. 
\begin{theorem}\label{th2.5}
Seja $S$ um semigrupo de classe $C^0$ em $X$. Então,
\[\lim\limits_{t\to +\infty}\frac{\log(\nl{S(t)})}{t}=\inf\limits_{t>0}\frac{\log(\nl{S(t)})}{t}=:\omega_0\]
e para cada $\omega>\omega_0$, existe {$M\geq 1$} tal que 
\begin{equation*}
{\nl{S(t)}\leq Me^{\omega t},\, \forall t\geq 0.}
\end{equation*}
\end{theorem} 


\begin{remark}
Quando $\omega_0<0$, então para $\omega=0$, temos que 
\[\nl{S(t)}\leq M,\ \forall t\geq 0. \]
Neste caso, dizemos que $S$ é um \dt{semigrupo uniformemente limitado}. Se, além disso, $M=1$, $S$ é dito \dt{semigrupo das contrações}. \index{semigrupo das contrações}
    
\end{remark}


\begin{lemma}\label{lem2.5}
    Seja $p:[0,+\infty)\longrightarrow \R$ uma função {subaditiva}, isto é, $p(t+s)\leq p(t)+p(s)$. Se $p$ é limitada superiormente em todo intervalo limitado, então $p(t)/t$ tem um limite quanto $t\to +\infty$ e 
    \[\lim\limits_{t\to +\infty}\frac{p(t)}{t}=\inf\limits_{t>0}\frac{p(t)}{t}.\]
Prova: Ver \cite[Lema 1.2.5]{gomes1985semigrupos} 
\end{lemma}


\begin{proof}[Prova do Teorema \ref{th2.5}]
Primeiramente, vejamos que $p(t)=\log\left(\nl{S(t)}\right)$ é subaditiva. De fato,
como $S(t)\in \mathcal{L}(X)$, temos que
\[\nl{S(t+s)}\leq \nl{S(t)S(s)}\leq \nl{S(t)}\nl{S(s)},\, \forall t,s\geq 0.\]
Assim, como a função $\log$ é crescente, temos que
\begin{align*}
p(t+s)& =\log\left(\nl{S(t+s)}\right)=\log\left(\nl{S(t)}\nl{S(s)}\right)\\
& \leq\log\left(\nl{S(t)}\right)+\log\left(\nl{S(s)}\right)\\
& \leq p(t)+p(s).
\end{align*}

A fim de aplicarmos o lema anterior, resta mostrar que $p$ é limitada superimente em todo intervalo limitado. Com efeito, seja $(a,b)$ um intervalo limitado em $[0,+\infty)$. Em particular, $(a,b)\subset [0,b]$. Portanto, da desigualdade \eqref{des.wT}, temos que
\begin{equation*}
    \nl{S(t)}\leq Me^{\mu b},\, \forall t\in (a,b).
\end{equation*}
Isto é, $\nl{S(t)}$ é limitada superiormente em todo intervalo limitado. Como $\log$ é crescente, temos que  $p$ também o é. Logo, do Lema \ref{lem2.5}, temos que
\begin{equation*}
    \lim\limits_{t\to \infty}\frac{\log\left(\nl{S(t)}\right)}{t}
    =\inf\limits_{t>0}\frac{\log\left(\nl{S(t)}\right)}{t}=:\omega_0.
\end{equation*}
Se $\omega>\omega_0$, tome $\ep=\omega-\omega_0$, pela definição de limite, existe $t_0>0$ tal que se $t>t_0$, então
\[\frac{\log\left(\nl{S(t)}\right)}{t}-\omega_0<\ep=\omega-\omega_0.\]
Donde, 
\[\frac{\log\left(\nl{S(t)}\right)}{t}<\omega\Rightarrow 
\nl{S(t)} \leq e^{\omega t},\, \forall t>t_0.\]
Por outro lado, da desigualdade \eqref{des.wT}, temos que
\[
\nl{S(t)} \leq Me^{\mu t_0}=:M_0,\, \forall t\in [0,t_0].
\]
E como $S(0)=\id$, então $M_0\geq 1$.

\noindent\textbf{1º caso:} $\omega\geq 0$.\\
Vimos que
\begin{equation*}
    \begin{cases}
    \log\left(\nl{S(t)}\right)\leq \log(M_0), & 0\leq t\leq t_0,\\
        \log\left(\nl{S(t)}\right)< t\omega, & t>t_0.        
    \end{cases}
\end{equation*}
Com isso, 
\begin{equation*}
 \log\left(\nl{S(t)}\right)\leq \max\{ \log(M_0),t\omega\}\leq \log(M_0)+t\omega ,\, \forall t\in [0,+\infty).
\end{equation*}
Donde,
\begin{equation*}
    \nl{S(t)}\leq e^{\log(M_0)+t\omega}\leq M_0e^{t\omega},\, \forall t\in [0,+\infty). 
\end{equation*}

\noindent\textbf{2º caso:} $\omega< 0$.\\
Neste caso, se $t>t_0$, como $-t_0\omega\geq 0$ e $\log(M_0)\geq 0$, então
\begin{equation*}
\log\left(\nl{S(t)}\right)\leq t\omega<\underbrace{\log(M_0)-t_0\omega}_{\geq 0}+t\omega, \, \forall t>t_0. 
\end{equation*}

Por outro lado, se $t\leq t_0$, então $t\omega-t_0\omega\geq 0$, daí,
\begin{equation*}
\log\left(\nl{S(t)}\right)\leq \log(M_0)<\log(M_0)\underbrace{-t_0\omega+t\omega}_{\geq 0}, \, \forall 0\leq t\leq t_0. 
\end{equation*}

Em resumo,
\begin{equation*}
\log\left(\nl{S(t)}\right)\leq \log(M_0){-t_0\omega+t\omega}, \, \forall t\geq 0. 
\end{equation*}
Consequentemente,
\begin{equation*}
    \nl{S(t)}\leq \underbrace{M_0e^{-t_0\omega}}_{M}e^{t\omega}=Me^{t\omega}, \forall t\geq 0.
\end{equation*}
\end{proof}

\begin{definition}
    Seja $S$ um semigrupo de classe $C^0$ em $X$. O \dt{gerador infinitesimal}\index{Gerador Infinitesimal} de $S$ é o operador $A:D(A)\subset X\longrightarrow X$ definido por 
\[D(A)=\left\{x\in X;\ \lim\limits_{h\to 0^+} \frac{S(h)x-x}{h} \text{ existe}\right\},\]
\[Ax:=\lim\limits_{h\to 0^+} \frac{S(h)x-x}{h},\, x\in D(A)\]
\end{definition}

\begin{proposition}
     $D(A)$ é um subespaço vetorial de $X$ e $A$ é um operador linear.
\end{proposition}
\begin{proof}
    Exercício.
\end{proof}


\begin{remark}
De acordo com a definição acima, todo semigrupo tem um gerador infinitesimal associado. A principal questão é a recíproca, isto é, {\color{blue} quando um operador $A:D(A)\subset X\longrightarrow X$ é o gerador infinitesimal de algum semigrupo?} O Teorema de Hille-Yosida nos dará as condições para responder a essa pergunta.
\end{remark}


A seguir, apresentamos algumas notações que serão adotadas.
\begin{enumerate}
 \item Dado $S$ é um semigrupo de classe $C^0$ em $X$, vamos designar por $A_h$ o operador linear limitado 
    \[A_hx=\frac{S(h)x-x}{h},\ \forall x\in X.\]
\item Escrevemos $S(t)=e^{tA}$ para dizer que $A$ é o gerador infinitesimal de um semigrupo de classe $C^0$ em $X$.
\item Escrevemos $A\in G(M,\omega)$ para exprimir que $A$ é o gerador infinitesimal de um semigrupo de classe $C^0$, que satisfaz a condição:
\[\|e^{tA}\|\leq {M}e^{{\omega} t}, \forall t\geq 0.\]
\end{enumerate}

\begin{theorem}
Seja $S(t)=e^{tA}$ em $X$. Dado $x\in D(A)$, então 
\[e^{(\cdot)A}x\in C^0([0,+\infty);D(A))\cap C^1([0,+\infty);X)\]
e
\[\frac{d}{dt}\left(e^{tA}x\right)=Ae^{tA}x=e^{tA}Ax.\]
\end{theorem}
\begin{proof}\noindent


\noindent\textbf{Afirmação 1:} Se $x\in D(A)$, então $S(t)x\in D(A)$ e $A(S(t)x)=S(t)Ax$.

Dado $x\in D(A)$, seja $y=S(t)x$. Primeiramente, vamos mostrar que
$\lim\limits_{h\to 0^+} A_hy=S(t)Ax$. Para isso, note que
\begin{align*}
    A_hy&=\frac{S(h)y-y}{h}=\frac{S(h)S(t)x-S(t)x}{h}=\frac{S(t+h)x-S(t)x}{h}\\
    &=\frac{S(t)S(h)x-S(t)x}{h}=S(t)\frac{S(h)x-x}{h}=S(t)A_hx.
\end{align*}
Como $x\in D(A)$, temos que $\lim\limits_{h\to^0}A_hx=Ax$. Além disso, $S(t)\in \mathcal{L}(X)$, então $S(t)$ é contínua. Portanto, 
\[
\lim\limits_{h\to 0^+} A_hy=\lim\limits_{h\to 0^+} S(t)A_hx
=S(t)\left(\lim\limits_{h\to 0^+} A_hx\right)=S(t)Ax.
\]

Neste caso, provamos que $S(t)x=y\in D(A)$ e que $Ay=S(t)Ax$, ou seja, $A(S(t)x)=S(t)Ax$.

\noindent\textbf{Afirmação 2:} $\frac{d}{dt}S(t)x=A(S(t)x)=S(t)Ax,\ \forall x\in D(A).$

Primeiramente, note que
\begin{equation*}
\begin{split}
    A(S(t)x)& =\lim\limits_{h\to 0^+} A_h(S(t)x) =\lim_{h \to 0^+} \frac{S(h)S(t)x - S(t)x}{h}\\
    &= \lim_{h \to 0^+} \frac{S(t+h)x - S(t)x}{h} = \frac{d^+}{dt} S(t)x, \quad \forall x \in D(A).
\end{split}
\end{equation*}
Com isso, temos que 
    \begin{equation}\label{d+}
\frac{d}{dt}^+ S(t)x = A(S(t)x) = S(t)(Ax).    
\end{equation}

Agora, vamos calcular a derivada pela esquerda.
\[
\frac{d}{dt}^- S(t)x = \lim_{\delta \to 0^-} \frac{S(\overbrace{t+\delta}^{>0})x - S(t)x}{\delta}, \ \text{ para } -t < \delta < 0.
\]
Fazendo $\delta=-h$, temos que $0<h<t$ e 
\begin{equation}\label{eq2.3}
\begin{split}
    \frac{d}{dt}^- S(t)x 
    & = \lim_{h \to 0^+} \frac{S(t-h)x - S(t)x}{-h}=
    \lim_{h \to 0^+} \frac{S(t-h)x - S(t-h)S(h)x}{-h}\\
    &=  \lim_{h \to 0^-} S(t-h)\left(\frac{x - S(h)x}{-h}\right)
    =\lim_{h \to 0^+} S(t-h)A_hx\\
    &= \lim_{h \to 0^+} \Big(S(t-h)(A_hx-Ax)+S(t-h)Ax\Big).
\end{split}
\end{equation}
Do Corolário \eqref{Scontinua}, temos que $f(h)=S(t-h)Ax$ é contínua em $[0,t)$, portanto
\begin{equation}\label{eq2.4}
\lim_{h \to 0^+} S(t-h)Ax=\lim_{h \to 0^+}f(h)=f(0)=S(t)Ax.
\end{equation}
Por outro lado, do Teorema \ref{th2.5}, temos que
\[\nl{S(t-h)}\leq Me^{\omega(t-h)}\leq Me^{\omega t},\ \forall h\in [0,t),\]
donde,
\begin{equation}\label{eq2.5}
\begin{split}
 \lim_{h \to 0^+}\|S(t-h)(A_hx-Ax)\|
 & \leq \lim_{h \to 0^+}\nl{S(t-h)}\|A_hx-Ax\|\\
 & \leq \lim_{h \to 0^+}Me^{\omega t}\|A_hx-Ax\|=0.
\end{split}
\end{equation}
Com isso, \eqref{eq2.3}, \eqref{eq2.4} e \eqref{eq2.5} implicam que 
\begin{equation}\label{d-}
\frac{d}{dt}^-S(t)x=S(t)Ax=A(S(t)x),\ \forall x\in D(A).    
\end{equation}
Portanto, de \eqref{d+} e \eqref{d-}, temos que
\begin{equation*}
\frac{d}{dt}S(t)x=S(t)Ax=A(S(t)x),\ \forall x\in D(A).    
\end{equation*}
\end{proof}


\begin{example} \textbf{Existência e Unicidade de um PVI}

Seja $S(t)=e^{tA}$ em $X$. Se {$x_0\in D(A)$}, então $x(t)=e^{tA}x_0$ define {uma única solução} do PVI
\[
\begin{cases}
    x'(t)=Ax,\ t\in [0,+\infty)\\
    x(0)=x_0.
\end{cases}
\]
Se {$x_0\not\in D(A)$} em $X$, então $x(t)=e^{tA}x_0$ {não é diferenciável}. Neste caso, dizemos que $x=x(t)$ é uma \dt{solução branda} (\dt{mild solution}, em inglês.)\index{Solução!branda} \index{Solução!fraca} do PVI.
\end{example}

\begin{xca}
Seja $S(t)=e^{tA}$ em $X$. Se $x\in D(A)$ mostre que
\[e^{tA}x-e^{sA}x=\int_s^t Ae^{\tau A}x\,d\tau =\int_s^t e^{\tau A}Ax\,d\tau\] 
\end{xca}

\begin{proposition}
Se $S(t)=e^{tA}$ em $X$, então para todo $x\in X$, 
\[\int_0^t e^{sA}x\,ds\in D(A)\ \text{ e }\ A\left(\int_0^t e^{sA}x\,ds\right)=e^{tA}x-x.\]
\end{proposition}
\begin{proof}
    {\color{red} ESCREVER}
\end{proof}

\begin{proposition}
    Se $S(t)=e^{tA}$ em $X$, então $A$ é fechado e seu domínio é denso em $X$.
\end{proposition}
\begin{proof}
    {\color{red} ESCREVER}
\end{proof}

\begin{remark}
    Esta última proposição nos dá uma condição necessária para que um operador $A$ seja o gerador infinitesimal de um semigrupo.
\end{remark}

\begin{proposition}[Unicidade]
    Se $S_1(t)=e^{tA}$ e $S_2(t)=e^{tA}$ em $X$, então $S_1=S_2$.
\end{proposition}
\begin{proof}
    {\color{red} ESCREVER}
\end{proof}

\begin{definition}
    Seja $S(t)=e^{tA}$ em $X$. Defina $A^0=\operatorname{Id}$, $A^1=A$ e, supondo que $A^{n-1}$ esteja definido, vamos definir $A^n$ pondo:
\[D(A^n)=\{x\in D(A^{n-1});\ A^{n-1}x\in D(A)\},\]
\[A^nx=A(A^{n-1}x), \ \forall x\in D(A^{n}).\]
\end{definition}

\begin{proposition}
Seja $S(t)=e^{tA}$ em $X$. Temos:

\begin{enumerate}[(i)]
 \item $\mathcal{D}\left(A^n\right)$ é um subespaço de $X$ e $A^n$ é um operador linear de $X$;

 \item Se $x \in \mathcal{D}\left(A^n\right)$, então $S(t) x \in \mathcal{D}\left(A^n\right), \forall t \geq 0 \mathrm{e}$
\[
\frac{d^n}{d t^n} S(t) x=A^n S(t) x=S(t) A^n x, \forall n \in \mathbb{N}
\]
\item É válida a fórmula de Taylor: se $x \in \mathcal{D}\left(A^n\right)$, então
\[
S(t) x=\sum_{k=0}^{n-1} \frac{(t-a)^k}{k!} A^k S(a) x+\frac{1}{(n-1)!} \int_a^t(t-\tau)^{n-1} A^n S(\tau) x d \tau
\]
\item  $(S(t)-I)^n x=\int_0^t \cdots \int_0^t S\left(\tau_1+\cdots+\tau_n\right) A^n x\, d \tau_1 \cdots d \tau_n, \forall x \in \mathcal{D}\left(A^n\right)$;
\item $\displaystyle\bigcap_n \mathcal{D}\left(A^n\right)$ é denso em $X$.
\end{enumerate}
\end{proposition}

\begin{proof}
    {\color{red} ESCREVER}
\end{proof}
\section{Teorema de Hille-Yosida}



\begin{definition}
Seja $X$ um espaço de Banach e $A:D(A)\subset X\longrightarrow X$ um operador linear fechado. Chamamos de \dt{resolvente de $A$}, \index{Resolvente!conjunto} o seguinte conjunto
\[\rho(A)=\{\lambda \in \mathbb{C};\ \lambda \id -A \text{ é bijetor }\}\]
%O conjunto \dt{resolvente de $A$}, \index{Resolvente!conjunto} denotado por $\rho(A)$, é o conjunto de todos os $\lambda\in \mathbb{C}$ tais que:
%\begin{itemize}
%    \item $\lambda \id -A$ é bijetor;
%    \item $R(\lambda,A):=(\lambda \id -A)^{-1}: X\to $, chamado de \dt{operador resolvente} \index{Resolvente!operador} de $A$, é limitado.
%\end{itemize}
O conjunto $\sigma(A)=\mathbb{C}\setminus\rho(A)$ é dito \dt{espectro} \index{Espectro} de $A$. Para cada $\lambda\in \rho(A)$, definimos o  \dt{operador resolvente} \index{Resolvente!operador} por
\[R(\lambda,A):=(\lambda \id -A)^{-1}: X\longrightarrow D(A).\]
O qual, pelo Teorema do Gráfico fechado, é um operador linear limitado.


Note que, se $\lambda \in \rho(A)$, então para todo $x\in X$ 
\begin{equation*}
(\lambda \id -A)R(\lambda,A)x=x\Rightarrow \lambda \rla x-x=A\rla x.
\end{equation*}
Analogamente, se $x\in D(A)$, então
\begin{equation*}
R(\lambda,A)(\lambda \id -A)x=x\Rightarrow \lambda \rla x-x=\rla Ax.
\end{equation*}
Destas duas equações, podemos concluir que
\begin{equation}\label{RAAR}
A\rla x=\lambda \rla x-x=\rla Ax,\; \forall\, x\in D(A).
\end{equation}

\begin{xca}
Prove que se $\lambda, \mu\in \rho(A)$, então $\rla$ e $R(\mu,A)$ comutam.
\end{xca}

\end{definition}

\begin{theorem}[Hille-Yosida]\label{HY-contr}
Seja $X$ um espaço de Banach. Um operador linear $A:D(A)\subset X\longrightarrow X$ é o gerador infinitesimal de um 
\textbf{semigrupos de contrações} se, e somente se, 
\begin{enumerate}[$(i)$]
\item $A$ é fechado e densamente definido, i.e., $\overline{D(A)}=X$. 
\item $(0,+\infty)\subset \rho(A)$ e para todo $\lambda>0$
\begin{equation*}
\nl{R(\lambda,A)}\leq \frac{1}{\lambda}.
\end{equation*}
\end{enumerate}
\end{theorem}

A prova deste teorema será divida em  lemas.

\begin{lemma}[Condição Necessária]
Seja $A$ o gerador infinitesimal de um semigrupo de contrações. Então $A$ é fechado e densamente definido. Além disso, $(0,+\infty)\subset \rho(A)$, 
\[R(\lambda,A)x=\displaystyle\int_0^\infty e^{-\lambda \xi}S(\xi)x\, d\xi,\] 
e $\nl{R(\lambda, A)}\leq \displaystyle\frac{1}{\lambda}$.
\end{lemma}
\begin{proof} Da Proposição \ref{Afd}, temos que $A$ é fechado e densamente definido.

Dados $x\in X$ e $\lambda>0$ defina
\begin{equation*}
L_\lambda x =\int_0^\infty e^{-\lambda \xi}S(\xi)x\, d\xi.
\end{equation*}
Primeiramente, vejamos que $L_\lambda x$ está bem definida. 
Para isso, vamos usar o Teste de Weierstrass
(Proposição \ref{TWeierstrass}). Do Corolário \ref{Scontinua}, temos que a função
\begin{equation*}
(t,\lambda)\in [0,+\infty)\times (0,+\infty)\longmapsto e^{-\lambda t}S(t)x\in  X,
\end{equation*}
é contínua em $t$ para cada $\lambda\in (0,+\infty)$. 

Como $S$ é um semigrupo de contrações, temos que
\begin{equation*}
\|e^{-\lambda t}S(t)x\| =e^{-\lambda t}\|S(t)x\|\leq e^{-\lambda t}\|x\|
\leq e^{-t}\|x\| =:M(t), 
\; \forall\ (t,\lambda)\in [0,+\infty)\times (0,+\infty)
\end{equation*}
Além disso, 
\begin{align*}
 \int_0^\infty M(t)\, dt& =\int_0^\infty e^{-t}\|x\|\, dt
=\|x\| \int_0^\infty e^{- t}\, dt 
=\|x\|\lim_{b\to +\infty} \int_0^b e^{-t}\, dt\\
& =\|x\|\lim_{b\to +\infty}-e^{-t}\Big\vert_{t=0}^{t=b} =\|x\|\lim_{b\to +\infty}\left(1-e^{-b}\right)={\|x\|}<+\infty.
\end{align*}
Portanto, pelo Teste de Weierstrass, $L_\lambda$ está bem definida. Claro que $L_\lambda$ é linear e além disso, procedendo como acima, 
\begin{align*}
\left\|L_\lambda x\right\|& \leq \int_0^\infty\|e^{-\lambda t}S(t)x\|\,dt\leq
\int_0^\infty e^{-\lambda t}\|x\|\,dt\\
& = \frac{\|x\|}{\lambda}\lim_{b\to +\infty}-e^{-\lambda t}\Big\vert_{t=0}^{t=b}
=\frac{\|x\|}{\lambda},
\end{align*}
isto é, $L_\lambda$ define um operador linear limitado em $X$ e 
$\nl{L_\lambda}\leq \frac{1}{\lambda}$.

Resta mostrar que  $L_\lambda= R(\lambda,A)$, isto é, devemos mostrar que
\begin{enumerate}
\item $\forall\, x\in X,\ L_\lambda x\in D(A)$ e $(\lambda \id-A)L_\lambda x=x$, i.e., $A\left(L_\lambda x\right)=L_\lambda x -x$.
\item $\forall\, x\in D(A),$ $L_\lambda(\lambda \id-A)x=x$, i.e., $L_\lambda(Ax)=\lambda L_\lambda x-x$.
\end{enumerate}

De fato, dado $x\in X$, seja $h>0$, como $A_h\in \mathcal{L}$, temos que
\begin{align*}
A_h(L_\lambda x)& =A_h\left(\int_0^\infty e^{-\lambda \xi}S(\xi)x\, d\xi\right)
=\frac{1}{h}\int_0^\infty (S(h)-\id) e^{-\lambda \xi}S(\xi)x\, d\xi\\
&= \frac{1}{h}\int_0^\infty e^{-\lambda \xi}S(\xi+h)x\, d\xi- 
\frac{1}{h}\int_0^\infty e^{-\lambda \xi}S(\xi)x\, d\xi\\
\shortintertext{(Fazendo a mudança $s=\xi+h$ na primeira integral e trocando $\xi$ por $s$ na segunda)}
& = \frac{1}{h}\int_h^\infty e^{-\lambda (s-h)}S(s)x\, ds- 
\frac{1}{h}\int_0^\infty e^{-\lambda s}S(s)x\, ds\\
& = \frac{e^{\lambda h}}{h}\int_h^\infty e^{-\lambda s}S(s)x\, ds- 
\frac{1}{h}\int_0^\infty e^{-\lambda s}S(s)x\, ds\\
& = \frac{e^{\lambda h}}{h}\left(\int_0^\infty e^{-\lambda s}S(s)x\, ds
-\int_0^h e^{-\lambda s}S(s)x\, ds\right)- 
\frac{1}{h}\int_0^\infty e^{-\lambda s}S(s)x\, ds\\
& = \frac{e^{\lambda h}-1}{h}\int_0^\infty e^{-\lambda s}S(s)x\, ds
-\frac{e^{\lambda h}}{h}\int_0^h e^{-\lambda s}S(s)x\, ds\\
&= \frac{e^{\lambda h}-1}{h}L_\lambda x
-\frac{e^{\lambda h}}{h}\int_0^h e^{-\lambda s}S(s)x\, ds
\end{align*}
Aplicando-se o limite quando $h\to0^+$, obtemos
\begin{align*}
A\left(L_\lambda x\right)&=\lim_{h\to0^+}A_h(L_\lambda x) =\lim_{h\to0^+}\left(\frac{e^{\lambda h}-1}{h}
L_\lambda x-e^{\lambda h}\frac{1}{h}\int_0^h
e^{-\lambda s}S(s)x\, ds\right)\\
& =\frac{d}{dh}e^{\lambda h}\Big\vert_{h=0}
L_\lambda x-\lim_{h\to0^+}
\bigg(e^{\lambda h}\underbrace{\frac{1}{h}\int_0^h
\overbrace{e^{-\lambda s}S(s)x}^{\text{contínua}}\,
 ds}_{\substack{\downarrow \\ x}}\bigg)\\[-1em]
& =\lambda
L_\lambda x-x,
\end{align*}
o que prova o o item 1. Agora vamos provar o item 2. Dado $x\in D(A)$, 
\begin{align*}
L_\lambda(A)x& =\int_0^{\infty}e^{-\lambda
 \xi}S(\xi)Ax\,d\xi=\int_0^{\infty}e^{-\lambda
 \xi}\frac{d}{d\xi}S(\xi)x\,d\xi\\
& =\int_0^{\infty}\frac{d}{d\xi}\left(e^{-\lambda
 \xi}S(\xi)x\right)+\lambda e^{-\lambda
 \xi}S(\xi)x \,d\xi\\
& =\lim_{b\to +\infty}\left(e^{-\lambda
 \xi}S(\xi)x\big\vert_{\xi=0}^{\xi=b}\right)+\lambda \int_0^{\infty}e^{-\lambda
 \xi}S(\xi)x \,d\xi\\
& =\lim_{b\to +\infty}\left(e^{-\lambda
 b}S(b)x\right)-x+\lambda L_\lambda x\\
\shortintertext{(como $S$ é uma contração, $\|e^{-\lambda b}S(b)x\|\leq e^{-\lambda b}\|x\|\to 0$, quando $b\to +\infty$)}
& =-x+\lambda L_\lambda x,
\end{align*}
como queríamos.
\end{proof}

\begin{lemma} Seja $A$ satisfazendo as condições $(i)$ e $(ii)$ do Teorema \ref{HY-contr}. Então
\begin{equation*}
\lim_{\lambda \to +\infty} \lambda\rla x=x, \; \forall\,x\in X.
\end{equation*}
\end{lemma}
\begin{proof}
Se $x\in D(A)$, da identidade \eqref{RAAR}, temos que
\begin{equation*}
\|\lambda \rla x-x\|=\|\rla Ax\|\leq \frac{1}{\lambda}\|Ax\|\to 0\; 
\text{ quando } \lambda\to +\infty.
\end{equation*}
Agora, \underline{dado $x\in X$}, \underline{para cada $\ep>0$}, como $D(A)$ é denso em $X$, tome $y\in D(A)$ tal que
\[\|y-x\|< \frac{\ep}{4}.\]
E, da convergência anterior, como $y\in D(A)$,\underline{ existe $\lambda_0>0$ tal que se
$\lambda >\lambda_0$}, então
\begin{equation*}
\|\lambda \rla y-y\|< \frac{\ep}{2}.
\end{equation*}
Com isso, 
\begin{align*}
\underline{\|\lambda \rla x-x\|}& \leq \|\lambda \rla x-\lambda\rla y\|
+\|\lambda \rla y-y\|+\|y-x\|\\
& < \lambda \|\rla (x-y)\|+ \frac{\ep}{2}+\frac{\ep}{4}\\
& \leq \|x-y\|+ \frac{3\ep}{4}\\
&< \frac{\ep}{4}+\frac{3\ep}{4}=\underline{\ep}\\
\end{align*}
\end{proof}

\begin{definition}
Para cada $\lambda \in (0,+\infty) \cap \rho(A)$, definimos o operador $A_\lambda: X\longrightarrow X$, chamado \dt{aproximação de Yosida}\index{Aproximação de Yosida} de $A$, por 
\begin{equation*}
A_\lambda :=\lambda A\rla =\lambda^2\rla -\lambda \id\in \mathcal{L}(X).
\end{equation*}
\end{definition}

\begin{lemma}
 Seja $A$ satisfazendo as condições $(i)$ e $(ii)$ do Teorema \ref{HY-contr}. 
Se $A_\lambda$ é a aproximação de Yosida do operador $A$, então
\begin{equation*}
\lim_{\lambda \to +\infty}A_\lambda x=Ax,\; \forall\, x\in D(A).
\end{equation*}
\end{lemma}
\begin{proof}
Se $x\in D(A)$, então, de \eqref{RAAR} e do Lema anterior, temos que
\begin{equation*}
\lim_{\lambda \to +\infty}A_\lambda x=\lim_{\lambda \to +\infty}\lambda\rla Ax=Ax
\end{equation*}
\end{proof}

\begin{lemma}
Seja $A$ satisfazendo as condições $(i)$ e $(ii)$ do Teorema \ref{HY-contr}. Se $A_\lambda$ é  aproximação de Yosida de $A$, então $A_\lambda$ é o gerador infinitesimal de um semigrupo uniformemente contínuo de contrações $e^{tA_\lambda}$. Além disso, para cada $x\in X$, $\lambda,\mu>0$ temos que
\begin{equation*}
\left\|e^{tA_\lambda}x-e^{tA_\mu}x\right\|\leq t\|A_\lambda x-A_\mu x\|.
\end{equation*}
\end{lemma}
\begin{proof}
Como $A_\lambda \in \mathcal{L}(X)$, então $A_\lambda$ é o gerador infinitesimal do semigrupo uniformemente contínuo $e^{tA_\lambda}$. Além disso,
\begin{multline*}
\left\|e^{tA_\lambda}\right\|
 =\left\|e^{t\lambda^2\rla-t\lambda \id}\right\|
\leq \left\|e^{t\lambda^2\rla}\right\|\left\|e^{-t\lambda \id} \right\|
=e^{-t\lambda }\left\|e^{t\lambda^2\rla}\right\|\\ 
\leq e^{-t\lambda }e^{t\lambda^2\left\|\rla\right\|}
 \leq e^{t\lambda }e^{-t\lambda }=1
\end{multline*}
portanto, um semigrupo de contrações.

Das definições de $A_\lambda$, $A_\mu$, $e^{tA_\lambda}$ e $e^{tA_\mu}$ 
comutam. Como $A_\lambda,  A_\mu\in\mathcal{L}(X)$, defina 
\[f(s)=e^{tsA_\lambda}e^{t(1-s)A_\mu}x=e^{tA_\mu}e^{st(A_\lambda-A_\mu)} \]
e note que
\begin{enumerate}
\item $f(1)=e^{tA_\lambda}x$, $f(0)=e^{tA_\mu}x$
\item $f'(s)=t(A_\lambda-A_\mu)e^{tsA_\lambda}e^{t(1-s)A_\mu}x$
\end{enumerate}
Consequentemente, pelo Teorema Fundamental do Cálculo,
\begin{align*}
\left\|e^{tA_\lambda}x-e^{tA_\mu}x \right\|
& =\left\|f(1)-f(0) \right\|=
\left\|\int_0^1 f'(s)\,ds \right\|\\
& =\left\|\int_0^1 
t(A_\lambda-A_\mu)e^{tsA_\lambda}e^{t(1-s)A_\mu}x\right\|\\
& \leq t\underbrace{\left\|e^{tsA_\lambda}e^{t(1-s)A_\mu}\right\|}_{\leq 1}   \left\|A_\lambda x-A_\mu x\right\|\\
& \leq t \left\|A_\lambda x-A_\mu x\right\|
\end{align*}
\end{proof}











\appendix
                                 
                            \chapter{Apêndice}

\section{Resultados Clássicos}


\begin{proposition}[Teste de Weierstrass]\label{TWeierstrass}
Seja $f:[a,+\infty)\times \Lambda\longrightarrow X$, $\Lambda$ 
$\overbrace{\text{um subconjunto aberto de }\mathbb{C}}^{\text{Não seria } \R\,?}$ 
 contínua em $t\in [a,+\infty)$ para cada $\lambda \in \Lambda$ Se existe $M:[a,+\infty)\longrightarrow \R$ contínua e positiva em $t\in [a,+\infty)$ tal que
\begin{enumerate}[$(i)$]
\item $\|f(t,\lambda)\|\leq M(t),\quad \forall (t,\lambda)\in [a,+\infty)\times \Lambda$,
\item $\displaystyle \int_a^\infty M(t)\, dt<+\infty$.
\end{enumerate} 
Então 
\[\int_a^\infty f(t,\lambda)\, dt\]
converge absolutamente para cada $\lambda $ pertecente ao conjunto $\Lambda$ e a convergência é uniforme nesse conjunto.
\end{proposition}
\begin{proof}
{\color{red}CITAR}
\end{proof}


A seguir, enunciamos dois resultados básicos da Análise Funcional que serão utilizados nas exposições deste minicurso. O primeiro deles é o Princípio da Limitação Uniforme (versão do Teorema de Banach-Steinhaus no contexto dos espaços normados) e o outro é o Teorema do Gráfico Fechado, ambos consequências do importante Lema de Baire (veja \cite{pombo1999introduccao}).

\begin{theorem}[Banach-Steinhaus]\label{th-BS}
	Sejam $X$ e $Y$ dois espaços normado, com $X$ completo, e consideremos uma família
    \[
    \displaystyle \mathcal F =\{T_i :X\longrightarrow Y; i\in I\} \subset \mathcal L (X,Y),
    \]
    não necessariamente enumerável, com a seguinte propriedade: para cada $x\in X$, temos
    \[
    \sup_{i\in I}\|T_ix\|_{Y} < \infty .
    \] 
    Então 
    \[
    \sup_{i\in I}\|T_i\|_{\mathcal{L}(X,Y)}<\infty.
    \]
    Em outras palavras, a limitação pontual da família $\mathcal F$ implica a sua limitação uniforme. 
    %Se $\{T_i\}_{i\in I}$  é uma família (não necessariamente enumerável) em $\mathcal{L}(X,Y)$ pontualmente limitada, então $\{T_i\}_{i\in I}$ é uniformemente limitada, isto é, se 
	% \[\sup_{i\in I}\|T_ix\|< \infty,\ \forall x\in X,\] 
%	 então
	% \[\sup_{i\in I}\|T_i\|_{\mathcal{L}(X,Y)}<\infty.\]
\end{theorem}

\begin{remark}
\begin{itemize}
\item[(a)] No Teorema \ref{th-BS}, a hipótese de que $X$ seja coompleto não pode ser removida. De fato, definindo 
\[
\displaystyle \varphi_n : x=(x_j)_{j=1}^{\infty} \in c_{00} \longmapsto nx_n \in \mathbb K ,
\]
para cada inteiro positivo $n$, fica estabelecida uma família pontualmente limitada que não é uniformemente limitada (relembre o Exemplo \ref{incompleto} e consulte \cite{pombo1999introduccao}).
\item[(b)] Ainda com as notações do Teorema de \ref{th-BS}, a limitação uniforme de $\mathcal F$, expressa por 
\[
\sup_{i\in I}\|T_i\|_{\mathcal{L}(X,Y)}<\infty ,
\]
é o mesmo que dizer que $\mathcal F$ é uma família equicontínua (de fato, basta adaptar a demonstração da Proposição \ref{continuous} para obter este fato).
\end{itemize}
\end{remark}

\begin{theorem}[Gráfico Fechado]\label{th-GF}
    Sejam $X$ e $Y$ dois espaços de Banach, e consideremos uma aplicação linear $T:X\longrightarrow Y$. Se
    \[
    \displaystyle G_T = \{ (x,y)\in X\times Y; y=Tx\}
    \]
    é um subespaço fechado de $X\times Y$, então $T$ é contínua.
\end{theorem}

\begin{remark}
    \begin{itemize}
    \item[(a)] Dentre as hipóteses do Teorema \ref{th-GF}, a completude de cada um dos dois espaços não pode ser removida (veja \cite{pombo1999introduccao}).
    \item[(b)] Conforme o enunciado do Teorema do \ref{th-GF}, $G_T$ é um subespaço vetorial de $X\times Y$ (não apenas um subconjunto do mesmo, como ocorre em geral para aplicações entre dois conjuntos dados). Isto é verdadeiro porque $T:X\longrightarrow Y$ é uma aplicação linear. 
    \item[(c)] É fácil ver que a aplicação linear $T:X\longrightarrow Y$ se exprime como a composição
    \[
    T=\pi_2 \circ (\pi_1 |_{G_T})^{-1},
    \]
    onde $\pi_1$ e $\pi_2$ são as projeções naturais de $X\times Y$ sobre $X$ e $Y$, respectivamente. A principal dificuldade na obtenção do Teorema \ref{th-GF} é demonstrar que $(\pi_1 |_{G_T})^{-1}: X\longrightarrow G_T$ é uma aplicação contínua (veja \cite{pombo1999introduccao}).
    \end{itemize}
\end{remark}

A seguir, apresentamos uma importante consideração sobre o operador resolvente mencionado na Definição \ref{res}.

\begin{remark}\label{reslim}
    Consideremos as notações da Definição \ref{res}. Para cada $\lambda \in \rho (A)$, vejamos que o operador linear $\rla :X\longrightarrow D(A) \subset X$ é limitado. Com efeito, como $X$ é um espaço de Banach, o Teorema do \ref{th-GF} garante que é suficiente verificarmos que o gráfico $G$ de $\rla$ é um subespaço fechado de $X\times X$. Nesse caso, consideremos uma sequência $(x_n ,\rla x_n )_{n=1}^{\infty}$ em $G$ convergindo a $(x,y)\in X\times X$. Isto significa que
    \begin{equation}\label{conv}
    \displaystyle x_n \to x \text{  e 
 } \rla x_n \to y \text{ em } X.
    \end{equation}
    Da Observação \ref{rl} (a), para todo $n\in \mathbb N$, temos
    \[
    \displaystyle A\rla x_n = \lambda \rla x_n -x_n 
    %\Longrightarrow x_n = \lambda \rla x_n - A\rla x_n \in D(A).
    \]
    %Portanto, pela Observação \ref{rl}, os operadores $A$ e $\rla$ comutam em cada termo da sequência $(x_n )_{n=1}^{\infty}$, o que implica
    %\[
    %\displaystyle x_n = \lambda \rla %x_n - \rla A x_n 
    %\]
    %para todo $n\in \mathbb N$. 
    Das convergências em \eqref{conv}, temos
    \[
    \displaystyle\lim_{n\to +\infty} A\rla x_n 
    %=\lim_{n\to +\infty} \rla Ax_n 
    = \lim_{n\to +\infty} [\lambda \rla x_n - x_n] = \lambda y-x \text{  em  } X.
    \]
    Como 
    \[
    \displaystyle y_n :=\rla x_n \to y \text{  e   } Ay_n  = A\rla x_n \to \lambda y-x
    \]
    em $X$, o fato de $A$ ser um operador fechado nos diz que $y\in D(A)$ e que $Ay = \lambda y-x$. Logo,
    \[
    \displaystyle (\lambda \id -A)y =x \Longrightarrow \rla x=y \Longrightarrow (x,y)\in G.
    \]
\end{remark}


\section{Espaços $\ell_p$}

\begin{example}\label{lp}
    Seja $p\in [1,+\infty )$ e definamos 
    \[
    \displaystyle \ell_p =\left\{ (x_n )_{n=1}^{\infty}; \sum_{n=1}^{\infty} |x_n|^p <+\infty \right\}
    \]
    
    Mostraremos que as operações 
    \[
    \bigg( (x_n )_{n=1}^{\infty},(y_n )_{n=1}^{\infty} \bigg) \in \ell_p \times \ell_p \longmapsto (x_n +y_n )_{n=1}^{\infty} \in \ell_p
    \]
    e
    \[
    \bigg( \lambda ,(x_n )_{n=1}^{\infty} \bigg) \in \mathbb K \times \ell_p \longmapsto (\lambda x_n )_{n=1}^{\infty} \in \ell_p
    \]
    estão bem definidas e que a aplicação 
    \[
    \displaystyle (x_n )_{n=1}^{\infty} \in \ell_p \longmapsto \bigg( \sum_{n=1}^{\infty} |x_n|^p \bigg)^{1/p} \in \R .
    \]
    Tais objetivos serão alcançados no Corolário \ref{dms}.
\end{example}

\begin{proposition}[Desigualdade de Hölder para somas] \label{dhsomas}
    Sejam $p,q \in (1,+\infty)$ tais que 
    \[
    \displaystyle \frac{1}{p} + \frac{1}{q} =1.
    \]
    Se $(x_1, \cdots , x_n) \in \mathbb K^n$ e $(y_1, \cdots , y_n) \in \mathbb K^n$, então 
    \[
    \displaystyle \sum_{j=1}^{n} |x_j y_j|\leq \bigg( \sum_{j=1}^{n} |x_j|^p \bigg)^{1/p} \bigg( \sum_{j=1}^{n} |y_j|^q \bigg)^{1/q} .
    \]
\end{proposition}
\begin{proof}
    Em virtude do Teorema do Valor Médio, não é difícil constatar que, para quaisquer $a,b\in [0,+\infty )$ e $\alpha \in [0,1]$, tem-se
    \begin{equation}\label{TVM}
    \displaystyle a^{\alpha} b^{1-\alpha} \leq \alpha a + (1-\alpha )b.
    \end{equation}
    Como claramente podemos supor que
    \[
    \sum_{j=1}^{n} |x_j|^p >0 \text{  e  } \sum_{j=1}^{n} |y_j|^p >0,
    \]
    consideremos 
    \[
    \displaystyle a_{m} = \frac{|x_{m}|^p}{\sum_{j=1}^{n} |x_j|^p} \text{  e  } b_{m} = \frac{|y_{m}|^q}{\sum_{j=1}^{n} |y_j|^p},
    \]
    para cada $m\in \{1,\cdots , n\}$. Logo, tomando $\alpha = \frac{1}{p}$ e aplicando \eqref{TVM}, temos
    \[
    \displaystyle a_{m}^{\alpha} b_{m}^{1-\alpha} \leq \alpha a_m + (1-\alpha )b_m ,
    \]
    isto é,
    \begin{equation}\label{ineq}
    \displaystyle \frac{|x_{m}|}{\bigg(\sum_{j=1}^{n} |x_j|^p \bigg)^{1/p}} \cdot \frac{|y_{m}|}{\bigg(\sum_{j=1}^{n} |y_j|^p \bigg)^{1/q}}
    \leq \frac{1}{p} \cdot \frac{|x_{m}|^p}{\sum_{j=1}^{n} |x_j|^p}+ \frac{1}{q} \cdot \frac{|y_{m}|^q}{\sum_{j=1}^{n} |y_j|^p},
    \end{equation}
    para cada $m\in \{ 1,\cdots , n\}$. Somando membro a membro as $n$ relações descritas em \eqref{ineq}, fica demonstrada a desigualdade de Hölder do enunciado.
\end{proof}

\begin{corollary}[Desigualdade de Hölder para séries]
    Sejam $p,q \in (1,+\infty)$ satisfazendo $\frac{1}{p} + \frac{1}{q} =1$. Se $(x_n)_{n=1}^{\infty} \in \ell _p$ e $(y_n)_{n=1}^{\infty} \in \ell _q$, então $(x_n y_n)_{n=1}^{\infty} \in \ell _1$ e vale
    \[
    \displaystyle \sum_{n=1}^{\infty} |x_n y_n|\leq \bigg( \sum_{n=1}^{\infty} |x_n|^p \bigg)^{1/p} \bigg( \sum_{n=1}^{\infty} |y_n|^q \bigg)^{1/q} .
    \]
\end{corollary}
\begin{proof}
Segue imediatamente da Proposição \ref{dhsomas}.
\end{proof}


\begin{proposition}[Desigualdade de Minkowski para somas]
Seja $p\in (1,+\infty)$. Se $(x_1, \cdots , x_n) \in \mathbb K^n$ e $(y_1, \cdots , y_n) \in \mathbb K^n$, então 
    \[
    \displaystyle \bigg( \sum_{j=1}^{n} |x_j + y_j|^p \bigg)^{1/p} \leq \bigg( \sum_{j=1}^{n} |x_j|^p \bigg)^{1/p} +\bigg( \sum_{j=1}^{n} |y_j|^p \bigg)^{1/p} .
    \]
\end{proposition}

\begin{proof}
Como claramente podemos supor que
\[
\displaystyle \sum_{j=1}^{n} |x_j + y_j|^p >0,
\]
a Proposição \ref{dhsomas} nos dá

\begin{align}
\displaystyle \sum_{j=1}^{n} |x_j + y_j|^p
&=\sum_{j=1}^{n} |x_j + y_j|^{p-1} |x_j + y_j| \nonumber \\
&\leq \sum_{j=1}^{n} |x_j + y_j|^{p-1} |x_j|
+\sum_{j=1}^{n} |x_j + y_j|^{p-1} |y_j| \nonumber \\
&\leq \bigg( \sum_{j=1}^{n} |x_j + y_j|^{(p-1)q} \bigg)^{1/q} \bigg( \sum_{j=1}^{n} |x_j|^p \bigg)^{1/p} \nonumber\\
&\hspace{1cm}+\bigg( \sum_{j=1}^{n} |x_j + y_j|^{(p-1)q} \bigg)^{1/q} \bigg( \sum_{j=1}^{n} |y_j|^p \bigg)^{1/p} \nonumber \\
&=\bigg( \sum_{j=1}^{n} |x_j + y_j|^{p} \bigg)^{1/q} \bigg[ \bigg( \sum_{j=1}^{n} |x_j|^p \bigg)^{1/p} + \bigg( \sum_{j=1}^{n} |y_j|^p \bigg)^{1/p}\bigg] \label{ineqm},
\end{align}
onde $q=\frac{p}{p-1}$. Multiplicando os membros de \eqref{ineqm} por $\bigg( \sum_{j=1}^{n} |x_j + y_j|^p \bigg)^{-1/q}$, fica demonstrada a desigualdade de Minkowski do enunciado.
\end{proof}

\begin{corollary}[Desigualdade de Minkowski para séries]\label{dms}
    Seja $p\in (1,+\infty)$. Se $x=(x_n)_{n=1}^{\infty} \in \ell_p $ e $y=(y_n)_{n=1}^{\infty} \in \ell_q$, então $(x_n + y_n)_{n=1}^{\infty} \in \ell_p$ e vale
     \[
    \displaystyle \bigg( \sum_{n=1}^{\infty} |x_n + y_n|^p \bigg)^{1/p} \leq \bigg( \sum_{n=1}^{\infty} |x_n|^p \bigg)^{1/p} +\bigg( \sum_{n=1}^{\infty} |y_n|^p \bigg)^{1/p} .
    \]
    Em outras palavras, $x+y \in \ell_p$, ou seja, a adição em $\ell_p$ apresentada no Exemplo \ref{lp} encontra-se bem definida e, além disso, vale a desigualdade triangular
    \[
    \displaystyle \|x+y\|_p \leq \|x\|_p +\|y\|_p.
    \]
    Assim, sendo imediata a boa definição da multiplicação por escalar em $\ell_p$, também declarada no Exemplo \ref{lp}, resulta que $(\ell_p , \|\cdot\|_p)$ é um espaço vetorial normado.
\end{corollary}    


\begin{example}\label{ex24}
Denotemos por
\begin{itemize}
\item $\textbf{c}$ o conjunto de todas as sequências convergentes em $\mathbb K$;
\item $\textbf{c}_0$ o conjunto de todas as sequências convergentes em $\mathbb K$ convergindo para zero;
\item $\textbf{c}_{00}$ o conjunto de todas as sequências $(x_n)_{n=1}^{\infty}$ em $\mathbb K$ com a seguinte propriedade: existe um inteiro positivo $n_0$ tal que $x_n=0$ para todo $n\geq n_0$.
\end{itemize}
Claramente, 
\[
\textbf{c}_{00} \subset \textbf{c}_{0} \subset \textbf{c} \subset \ell_{\infty}.
\]
Além disso, não é difícil constatar que $\textbf{c}$ é um subespaço vetorial fechado de $\ell_{\infty}$.
\end{example}

\begin{example}
    Para cada espaço métrico $M$, $\mathcal C _b (M)$, apresentado no \ref{ex23} é um espaço de Banach, uma vez que é um subespaço fechado de $\mathcal B (M)$ (veja a Observação \ref{subespaço}).
\end{example}

\begin{example}
O subespaço vetorial $\textbf{c}$ de $\ell_{\infty}$, mencionado no Exemplo \ref{ex24}, é um espaço de Banach, em virtude da Proposição \ref{continuous} e da Observação \ref{subespaço}.
\end{example}

\begin{example}\label{incompleto}
O subespaço de $\textbf{c}_{00}$ de $\ell_{\infty}$, mencionado no Exemplo \ref{ex24}, NÃO é um espaço de Banach. De fato, para cada inteiro $n \geq 1$, consideremos
\[
\displaystyle x_n =\bigg( 1,\frac{1}{2}, \frac{1}{3}, \cdots , \frac{1}{n}, 0, 0, 0,\cdots \bigg) \in \textbf{c}_{00}.
\]
É fácil ver que $(x_n)_{n=1}^{\infty}$ é uma sequência em $\textbf{c} _{00}$ que converge a 
\[
\displaystyle x =\bigg( 1,\frac{1}{2}, \frac{1}{3}, \cdots , \frac{1}{n}, \cdots \bigg) \in \ell_{\infty} \setminus\textbf{c}_{00}.
\]
Isto significa que $(x_n)_{n=1}^{\infty}$ é uma sequência de Cauchy em $\textbf{c}_{00}$ que NÃO CONVERGE em $\textbf{c}_{00}$.
\end{example}

\section{Funcionais e operadores lineares ilimitados}


\begin{proposition}\label{dual-topVSalg}
Se é $X$ um espaço normado de dimensão infinita, então $X'\neq X^{\ast}$.
\end{proposition}
\begin{proof}
De fato, em virtude do Lema de Zorn, sabemos que $X$ possui uma base (de Hamel)
\[
\displaystyle \mathcal B = \{ x_i; i\in I\},
\]
onde $I$ é um conjunto infinito. Podemos supor que $\| x_i \|_X =1$ para todo $i\in I$. Agora, seja 
\[
\displaystyle J = \{ i_1, i_2, \cdots , i_k, \cdots \}
\]
um subconjunto infinito e enumerável de $I$. Logo, existe um único funcional linear $\varphi : X \longrightarrow \mathbb K$ tal que 
\begin{itemize}
\item $\varphi (x_{i_k})=k$ para cada $k\in \mathbb N$;
\item $\varphi (x_i)=0$ se $i\in I\setminus J$.
\end{itemize}
Um vez que 
\[
\displaystyle \mathbb N \subset \left\{ \frac{|\varphi (x)|}{\|x\|_{X}}; x\in X\setminus \{0\}\right\},
\]
não existe $C>0$ de modo que valha $|\varphi (x)|\leq \|x\|_X$ para todo $x\in X$. Pela Proposição \ref{continuous}, $\varphi$ não é contínua, isto é, $\varphi \in X^{\ast} \setminus X'$.
\end{proof}

\begin{remark}\label{opilimitado}
Dados dois espaços normados $X$ e $Y$, com $X$ de dimensão infinita, sempre podemos definir uma aplicação linear descontínua $T:X\longrightarrow Y$. Realmente, consideremos $\mathcal B$ e $I$ como na Proposição \ref{dual-topVSalg}, e seja $y\in Y$ com $\|y\|_Y =1$. Nesse caso, basta tomar $T:X\longrightarrow Y$ como sendo a única aplicação linear satisfazendo
\begin{itemize}
\item $T(x_{i_k})=ky$ para cada $k\in \mathbb N$;
\item $T(x_i)=0$ se $i\in I\setminus J$.
\end{itemize}
\end{remark}
 

\backmatter
\bibliographystyle{acm}
\bibliography{bibliografia}
\addcontentsline{toc}{chapter}{Bibliografia}

\printindex
\end{document}


%%%%%%%%%%%%%%%%%%%%%%%%%%%%%%%%%%%%%%%%%%%%%%%%%
